\documentclass[12pt,a4paper]{article}
\usepackage{fontspec}
\setmainfont{Times New Roman}[Path=C:/Windows/Fonts/,UprightFont=times.ttf,BoldFont=timesbd.ttf,ItalicFont=timesi.ttf,BoldItalicFont=timesbi.ttf]
\usepackage{amsmath,amssymb,amsfonts}
\usepackage[left=1.5cm,right=1.5cm,top=1.5cm,bottom=1.7cm,headheight=17pt]{geometry}
\usepackage{enumitem,array,fancyhdr,booktabs}

\pagestyle{fancy}
\fancyhf{}
\lhead{\small\textbf{HƯỚNG DẪN GIẢI CHI TIẾT -- ÔN TẬP TOÁN 12}}
\rhead{\small\textit{Biên soạn: Dương Gia Bảo}}
\rfoot{\small Trang \thepage}
\renewcommand{\headrulewidth}{0.4pt}
\fancypagestyle{firstpage}{\fancyhead{}\renewcommand{\headrulewidth}{0pt}}

\setlength{\parindent}{0pt}
\setlength{\parskip}{4pt}
\setlength{\emergencystretch}{1em}

\setlist[enumerate,1]{label=\textbf{Câu \arabic*:},leftmargin=*,itemsep=8pt,topsep=5pt}

\begin{document}
\thispagestyle{firstpage}

\noindent
\begin{minipage}[t]{0.36\textwidth}
    \centering
    {\fontsize{12}{14}\selectfont\textbf{BÀI TẬP VỀ NHÀ}}\\[4pt]
    \textit{Biên soạn: Dương Gia Bảo}\\[6pt]
    \setlength{\fboxrule}{0.8pt}
    \setlength{\fboxsep}{5pt}
    \fbox{\begin{minipage}{0.88\linewidth}
        \centering
        {\fontsize{12}{14}\selectfont\textbf{ĐÁP ÁN \& LỜI GIẢI}}
    \end{minipage}}
\end{minipage}\hfill
\begin{minipage}[t]{0.62\textwidth}
    \centering
    {\fontsize{13}{15}\selectfont\textbf{CHUYÊN ĐỀ ÔN TẬP SAU THI GIỮA KÌ I}}\\[3pt]
    {\fontsize{12}{14}\selectfont\textbf{THỐNG KÊ GHÉP NHÓM \& NGUYÊN HÀM -- TÍCH PHÂN}}\\[3pt]
    {\fontsize{11}{13}\selectfont\textit{(Chương trình GDPT 2018 -- Gồm 50 câu)}}
\end{minipage}

\par\vspace{4pt}
\hrule height 0.6pt
\vspace{8pt}

\begin{center}
    {\fontsize{13}{15}\selectfont\textbf{BẢNG ĐÁP ÁN NHANH 50 CÂU}}
\end{center}

\begin{center}
\renewcommand{\arraystretch}{1.2}
\begin{tabular}{|c|c||c|c||c|c||c|c||c|c|}
\hline
\textbf{Câu} & \textbf{Đ/A} & \textbf{Câu} & \textbf{Đ/A} & \textbf{Câu} & \textbf{Đ/A} & \textbf{Câu} & \textbf{Đ/A} & \textbf{Câu} & \textbf{Đ/A} \\ \hline
\textbf{1} & \textbf{A} & \textbf{11} & $x=10$; $6{,}55$ & \textbf{21} & \textbf{B} & \textbf{31} & $42$ & \textbf{41} & $4$ \\ \hline
\textbf{2} & \textbf{B} & \textbf{12} & $38{,}5$ & \textbf{22} & \textbf{A} & \textbf{32} & $5$ & \textbf{42} & $4$ \\ \hline
\textbf{3} & \textbf{A} & \textbf{13} & $56{,}43$ & \textbf{23} & \textbf{C} & \textbf{33} & $7$ & \textbf{43} & $50$ \\ \hline
\textbf{4} & \textbf{A} & \textbf{14} & $3$ & \textbf{24} & \textbf{A} & \textbf{34} & $-1$ & \textbf{44} & $35$ \\ \hline
\textbf{5} & \textbf{A} & \textbf{15} & $13{,}39$ & \textbf{25} & \textbf{B} & \textbf{35} & $2$ & \textbf{45} & $4{,}5$ \\ \hline
\textbf{6} & $25$ & \textbf{16} & \textbf{A} & \textbf{26} & \textbf{C} & \textbf{36} & $5$ & \textbf{46} & $16$ \\ \hline
\textbf{7} & $0{,}29$ & \textbf{17} & \textbf{B} & \textbf{27} & \textbf{D} & \textbf{37} & $11$ & \textbf{47} & $9$ \\ \hline
\textbf{8} & $3{,}06$ & \textbf{18} & \textbf{C} & \textbf{28} & \textbf{A} & \textbf{38} & $3$ & \textbf{48} & $37$ \\ \hline
\textbf{9} & $0{,}93$ & \textbf{19} & \textbf{A} & \textbf{29} & \textbf{B} & \textbf{39} & $12$ & \textbf{49} & $40$ \\ \hline
\textbf{10} & $0{,}96$ & \textbf{20} & \textbf{D} & \textbf{30} & \textbf{C} & \textbf{40} & $0{,}25$ & \textbf{50} & $35$ \\ \hline
\end{tabular}
\end{center}

\vspace{10pt}
\hrule height 0.6pt
\vspace{10pt}

\begin{center}
    {\fontsize{13}{15}\selectfont\textbf{HƯỚNG DẪN GIẢI CHI TIẾT}}
\end{center}

\subsection*{PHẦN 1. CHƯƠNG III: CÁC SỐ ĐẶC TRƯNG ĐO MỨC ĐỘ PHÂN TÁN}

\begin{enumerate}
\item \textbf{Đáp án A.}  
Theo định nghĩa trong SGK Toán 12, khoảng biến thiên của mẫu số liệu ghép nhóm là hiệu giữa đầu mút phải của nhóm cuối cùng và đầu mút trái của nhóm đầu tiên: $R = a_{k+1} - a_1$.

\item \textbf{Đáp án B.}  
Nhóm đầu tiên là $[12; 14)$ nên $a_1 = 12$. Nhóm cuối cùng là $[18; 20)$ nên $a_5 = 20$.  
Khoảng biến thiên là $R = 20 - 12 = 8$.

\item \textbf{Đáp án A.}  
Khoảng tứ phân vị của mẫu số liệu ghép nhóm được định nghĩa là hiệu giữa tứ phân vị thứ ba và tứ phân vị thứ nhất: $\Delta_Q = Q_3 - Q_1$.

\item \textbf{Đáp án A.}  
Theo công thức phương sai của mẫu số liệu ghép nhóm: $s^2 = \dfrac{1}{n}\sum\limits_{i=1}^k n_i (c_i - \bar{x})^2$.

\item \textbf{Đáp án A.}  
Vì điểm trung bình của An và Bình bằng nhau ($\bar{x} = 8{,}0$), bạn nào có độ lệch chuẩn nhỏ hơn thì các điểm số phân tán quanh giá trị trung bình ít hơn, nghĩa là kết quả học tập ổn định, đồng đều hơn. Vì $s_A = 0{,}65 < s_B = 1{,}20$ nên điểm của An đồng đều hơn Bình.

\item \textbf{Đáp số: $25$.}  
Nhóm đầu tiên là $[20; 25) \implies a_1 = 20$.  
Nhóm cuối cùng là $[40; 45) \implies a_6 = 45$.  
Khoảng biến thiên của mẫu số liệu ghép nhóm là:  
\[
R = 45 - 20 = 25\text{ (phút)}.
\]

\item \textbf{Đáp số: $0{,}29$.}  
Cỡ mẫu $n = 3 + 7 + 12 + 6 + 2 = 30$.  
- Tìm $Q_1$: Ta có $\dfrac{n}{4} = \dfrac{30}{4} = 7{,}5$.  
Tần số tích lũy: $cf_1 = 3 < 7{,}5$, $cf_2 = 3 + 7 = 10 \ge 7{,}5$.  
Do đó nhóm chứa $Q_1$ là $[1{,}2; 1{,}4)$. Ta có $u_2 = 1{,}2$; độ dài khoảng $h = 0{,}2$; tần số $n_2 = 7$; tần số tích lũy trước đó $C = 3$.  
\[
Q_1 = 1{,}2 + \dfrac{7{,}5 - 3}{7} \cdot 0{,}2 = 1{,}2 + \dfrac{4{,}5}{7} \cdot 0{,}2 = 1{,}2 + \dfrac{9}{70} = \dfrac{93}{70} \approx 1{,}3286\text{ kg}.
\]
- Tìm $Q_3$: Ta có $\dfrac{3n}{4} = \dfrac{90}{4} = 22{,}5$.  
Tần số tích lũy: $cf_2 = 10$, $cf_3 = 10 + 12 = 22 < 22{,}5$, $cf_4 = 22 + 6 = 28 \ge 22{,}5$.  
Do đó nhóm chứa $Q_3$ là $[1{,}6; 1{,}8)$. Ta có $u_4 = 1{,}6$; $h = 0{,}2$; $n_4 = 6$; $C' = 22$.  
\[
Q_3 = 1{,}6 + \dfrac{22{,}5 - 22}{6} \cdot 0{,}2 = 1{,}6 + \dfrac{0{,}5}{6} \cdot 0{,}2 = 1{,}6 + \dfrac{1}{60} = \dfrac{97}{60} \approx 1{,}6167\text{ kg}.
\]
Khoảng tứ phân vị: $\Delta_Q = Q_3 - Q_1 = \dfrac{97}{60} - \dfrac{93}{70} = \dfrac{121}{420} \approx 0{,}288 \implies \Delta_Q \approx 0{,}29\text{ kg}$.

\item \textbf{Đáp số: $3{,}06$.}  
Các giá trị đại diện của các nhóm: $c_1 = 1{,}5$; $c_2 = 2{,}5$; $c_3 = 3{,}5$; $c_4 = 4{,}5$.  
Cỡ mẫu $n = 8 + 15 + 18 + 9 = 50$.  
Số trung bình của mẫu số liệu ghép nhóm là:  
\[
\bar{x} = \dfrac{8 \cdot 1{,}5 + 15 \cdot 2{,}5 + 18 \cdot 3{,}5 + 9 \cdot 4{,}5}{50} = \dfrac{12 + 37{,}5 + 63 + 40{,}5}{50} = \dfrac{153}{50} = 3{,}06\text{ (giờ)}.
\]

\item \textbf{Đáp số: $0{,}93$.}  
Áp dụng công thức tính phương sai $s^2 = \dfrac{1}{n}\sum\limits_{i=1}^4 n_i c_i^2 - \bar{x}^2$:  
Ta tính:  
\[
\sum_{i=1}^4 n_i c_i^2 = 8 \cdot (1{,}5)^2 + 15 \cdot (2{,}5)^2 + 18 \cdot (3{,}5)^2 + 9 \cdot (4{,}5)^2
\]
\[
= 8 \cdot 2{,}25 + 15 \cdot 6{,}25 + 18 \cdot 12{,}25 + 9 \cdot 20{,}25 = 18 + 93{,}75 + 220{,}5 + 182{,}25 = 514{,}5.
\]
Phương sai:  
\[
s^2 = \dfrac{514{,}5}{50} - (3{,}06)^2 = 10{,}29 - 9{,}3636 = 0{,}9264 \approx 0{,}93.
\]

\item \textbf{Đáp số: $0{,}96$.}  
Độ lệch chuẩn là căn bậc hai của phương sai:  
\[
s = \sqrt{s^2} = \sqrt{0{,}9264} \approx 0{,}9625 \implies s \approx 0{,}96\text{ (giờ)}.
\]

\item \textbf{Đáp số: $x = 10$; $\bar{x} = 6{,}55$.}  
Tổng số học sinh $n = 40 \iff 4 + 10 + x + 12 + 4 = 40 \iff 30 + x = 40 \iff x = 10$.  
Giá trị đại diện các nhóm: $c_1 = 4{,}5; c_2 = 5{,}5; c_3 = 6{,}5; c_4 = 7{,}5; c_5 = 8{,}5$.  
Điểm trung bình:  
\[
\bar{x} = \dfrac{4 \cdot 4{,}5 + 10 \cdot 5{,}5 + 10 \cdot 6{,}5 + 12 \cdot 7{,}5 + 4 \cdot 8{,}5}{40} = \dfrac{18 + 55 + 65 + 90 + 34}{40} = \dfrac{262}{40} = 6{,}55.
\]

\item \textbf{Đáp số: $38{,}5$.}  
Cỡ mẫu $n = 30 \implies \dfrac{3n}{4} = \dfrac{90}{4} = 22{,}5$.  
Tần số tích lũy: $cf_1 = 5, cf_2 = 14, cf_3 = 24 \ge 22{,}5$.  
Nhóm chứa $Q_3$ là $[30; 40)$. Ta có $u_3 = 30$, độ dài $h = 10$, tần số $n_3 = 10$, tần số tích lũy trước đó $C = 14$.  
\[
Q_3 = 30 + \dfrac{22{,}5 - 14}{10} \cdot 10 = 30 + 8{,}5 = 38{,}5\text{ (triệu đồng)}.
\]

\item \textbf{Đáp số: $56{,}43$.}  
Cỡ mẫu $n = 40 \implies \dfrac{n}{4} = 10$.  
Tần số tích lũy: $cf_1 = 6 < 10, cf_2 = 6 + 14 = 20 \ge 10$.  
Nhóm chứa $Q_1$ là $[55; 60)$. Ta có $u_2 = 55, h = 5, n_2 = 14, C = 6$.  
\[
Q_1 = 55 + \dfrac{10 - 6}{14} \cdot 5 = 55 + \dfrac{20}{14} = 55 + \dfrac{10}{7} = \dfrac{395}{7} \approx 56{,}43\text{ kg}.
\]

\item \textbf{Đáp số: $3$.}  
Dây chuyền A: $s_A^2 = 36 \implies s_A = \sqrt{36} = 6\text{ g}$.  
Dây chuyền B: Cỡ mẫu $n = 5 + 15 + 20 + 10 = 50$.  
Giá trị đại diện các nhóm: $c_1 = 485, c_2 = 495, c_3 = 505, c_4 = 515$.  
Số trung bình của dây chuyền B:  
\[
\bar{x}_B = \dfrac{5 \cdot 485 + 15 \cdot 495 + 20 \cdot 505 + 10 \cdot 515}{50} = \dfrac{25100}{50} = 502\text{ g}.
\]
Phương sai của dây chuyền B:  
\[
s_B^2 = \dfrac{1}{50}\left[ 5(485-502)^2 + 15(495-502)^2 + 20(505-502)^2 + 10(515-502)^2 \right]
\]
\[
= \dfrac{1}{50}\left[ 5(-17)^2 + 15(-7)^2 + 20(3)^2 + 10(13)^2 \right] = \dfrac{1}{50}[1445 + 735 + 180 + 1690] = \dfrac{4050}{50} = 81.
\]
Độ lệch chuẩn của dây chuyền B: $s_B = \sqrt{81} = 9\text{ g}$.  
Hiệu độ lệch chuẩn: $s_B - s_A = 9 - 6 = 3\text{ g}$.

\item \textbf{Đáp số: $13{,}39$.}  
Cỡ mẫu $n = 40$.  
- Tìm $Q_1$: $\dfrac{n}{4} = 10$. Tần số tích lũy: $cf_1 = 6 < 10, cf_2 = 20 \ge 10 \implies$ Nhóm chứa $Q_1$ là $[10; 20)$.  
\[
Q_1 = 10 + \dfrac{10 - 6}{14} \cdot 10 = 10 + \dfrac{40}{14} = 10 + \dfrac{20}{7} = \dfrac{90}{7} \approx 12{,}86\text{ phút}.
\]
- Tìm $Q_3$: $\dfrac{3n}{4} = 30$. Tần số tích lũy: $cf_2 = 20 < 30, cf_3 = 20 + 16 = 36 \ge 30 \implies$ Nhóm chứa $Q_3$ là $[20; 30)$.  
\[
Q_3 = 20 + \dfrac{30 - 20}{16} \cdot 10 = 20 + \dfrac{100}{16} = 20 + 6{,}25 = 26{,}25\text{ phút}.
\]
Khoảng tứ phân vị: $\Delta_Q = Q_3 - Q_1 = 26{,}25 - \dfrac{90}{7} = \dfrac{375}{28} \approx 13{,}39\text{ phút}$.
\end{enumerate}

\subsection*{PHẦN 2. CHƯƠNG IV: NGUYÊN HÀM VÀ TÍCH PHÂN}

\begin{enumerate}[start=16]
\item \textbf{Đáp án A.}  
$\int (3x^2 - 2x + 5)\,\mathrm{d}x = 3 \cdot \dfrac{x^3}{3} - 2 \cdot \dfrac{x^2}{2} + 5x + C = x^3 - x^2 + 5x + C$.

\item \textbf{Đáp án B.}  
$\int (\sin x + 2\cos x)\,\mathrm{d}x = -\cos x + 2\sin x + C$.

\item \textbf{Đáp án C.}  
$\int (\mathrm{e}^{2x} - 3^x)\,\mathrm{d}x = \dfrac{1}{2}\mathrm{e}^{2x} - \dfrac{3^x}{\ln 3} + C$.

\item \textbf{Đáp án A.}  
$\int \dfrac{1}{2x+3}\,\mathrm{d}x = \dfrac{1}{2}\ln|2x+3| + C$. Do $x \in \left(-\dfrac{3}{2}; +\infty\right)$ nên $2x+3 > 0 \implies \dfrac{1}{2}\ln(2x+3) + C$.

\item \textbf{Đáp án D.}  
$F(x) = \int (4x^3 - 2x)\,\mathrm{d}x = x^4 - x^2 + C$.  
$F(1) = 1^4 - 1^2 + C = C = 5 \implies F(x) = x^4 - x^2 + 5$.  
Vậy $F(2) = 2^4 - 2^2 + 5 = 16 - 4 + 5 = 17$.

\item \textbf{Đáp án B.}  
Đặt $u = 3x - 1 \implies \mathrm{d}u = 3\,\mathrm{d}x \implies \mathrm{d}x = \dfrac{1}{3}\mathrm{d}u$.  
$\int f(3x-1)\,\mathrm{d}x = \dfrac{1}{3}\int f(u)\,\mathrm{d}u = \dfrac{1}{3}F(u) + C = \dfrac{1}{3}F(3x-1) + C$.

\item \textbf{Đáp án A.}  
Đặt $u = 2x+1 \implies \mathrm{d}u = 2\,\mathrm{d}x$; $\mathrm{d}v = \cos x\,\mathrm{d}x \implies v = \sin x$.  
Theo công thức từng phần: $\int u\,\mathrm{d}v = uv - \int v\,\mathrm{d}u = (2x+1)\sin x - 2\int \sin x\,\mathrm{d}x$.

\item \textbf{Đáp án C.}  
Theo tính chất cộng đoạn của tích phân:  
$\int_1^5 f(x)\,\mathrm{d}x = \int_1^3 f(x)\,\mathrm{d}x + \int_3^5 f(x)\,\mathrm{d}x = 4 + (-2) = 2$.

\item \textbf{Đáp án A.}  
$I = 2\int_0^2 f(x)\,\mathrm{d}x - \int_0^2 3x^2\,\mathrm{d}x = 2 \cdot 3 - \left[ x^3 \right]_0^2 = 6 - (8 - 0) = -2$.

\item \textbf{Đáp án B.}  
$I = \left[ x^3 + 2\mathrm{e}^x \right]_0^1 = (1 + 2\mathrm{e}) - (0 + 2\mathrm{e}^0) = 1 + 2\mathrm{e} - 2 = 2\mathrm{e} - 1$.

\item \textbf{Đáp án C.}  
Đặt $t = x^2+1 \implies \mathrm{d}t = 2x\,\mathrm{d}x$.  
Đổi cận: $x = 0 \implies t = 1$; $x = 1 \implies t = 2$.  
$I = \int_1^2 t^3\,\mathrm{d}t = \left[ \dfrac{t^4}{4} \right]_1^2 = \dfrac{16 - 1}{4} = \dfrac{15}{4}$.

\item \textbf{Đáp án D.}  
Đặt $u = x \implies \mathrm{d}u = \mathrm{d}x$; $\mathrm{d}v = \cos x\,\mathrm{d}x \implies v = \sin x$.  
$I = [x\sin x]_0^{\frac{\pi}{2}} - \int_0^{\frac{\pi}{2}} \sin x\,\mathrm{d}x = \left(\dfrac{\pi}{2} \cdot 1 - 0\right) - [-\cos x]_0^{\frac{\pi}{2}} = \dfrac{\pi}{2} + (\cos \frac{\pi}{2} - \cos 0) = \dfrac{\pi}{2} - 1$.

\item \textbf{Đáp án A.}  
Theo định nghĩa hình học của tích phân, diện tích hình phẳng giới hạn bởi đồ thị hàm số $y = f(x)$, trục hoành và hai đường thẳng $x = a, x = b$ ($a < b$) là $S = \int_a^b |f(x)|\,\mathrm{d}x$.

\item \textbf{Đáp án B.}  
Thể tích khối tròn xoay là:  
\[
V = \pi \int_1^4 (\sqrt{x})^2\,\mathrm{d}x = \pi \int_1^4 x\,\mathrm{d}x = \pi \left[ \dfrac{x^2}{2} \right]_1^4 = \pi \left(\dfrac{16}{2} - \dfrac{1}{2}\right) = \dfrac{15\pi}{2}.
\]

\item \textbf{Đáp án C.}  
Quãng đường đi được là tích phân của vận tốc từ $t = 0$ đến $t = 3$:  
$s = \int_0^3 (3t^2 + 2t)\,\mathrm{d}t = \left[ t^3 + t^2 \right]_0^3 = (27 + 9) - 0 = 36\text{ m}$.

\item \textbf{Đáp số: $42$.}  
$f(x) = \int f'(x)\,\mathrm{d}x = \int (6x^2 - 4x + 1)\,\mathrm{d}x = 2x^3 - 2x^2 + x + C$.  
Vì $f(1) = 4 \iff 2(1)^3 - 2(1)^2 + 1 + C = 4 \iff 1 + C = 4 \iff C = 3$.  
Do đó $f(x) = 2x^3 - 2x^2 + x + 3$.  
Giá trị $f(3) = 2(3)^3 - 2(3)^2 + 3 + 3 = 2(27) - 2(9) + 6 = 54 - 18 + 6 = 42$.

\item \textbf{Đáp số: $5$.}  
Biến đổi hàm số: $f(x) = \dfrac{2x+3}{x+1} = \dfrac{2(x+1) + 1}{x+1} = 2 + \dfrac{1}{x+1}$.  
Nguyên hàm trên $(-1; +\infty)$ là $F(x) = 2x + \ln(x+1) + C$.  
$F(0) = 2(0) + \ln(1) + C = 2 \iff C = 2$.  
Suy ra $F(x) = 2x + \ln(x+1) + 2$.  
Khi đó $F(1) = 2(1) + \ln(2) + 2 = 4 + \ln 2 \implies a = 4, b = 1 \implies T = a + b = 4 + 1 = 5$.

\item \textbf{Đáp số: $7$.}  
$F(x) = \int \dfrac{x}{\sqrt{x^2+9}}\,\mathrm{d}x$. Đặt $u = \sqrt{x^2+9} \implies u^2 = x^2+9 \implies u\,\mathrm{d}u = x\,\mathrm{d}x$.  
$F(x) = \int \dfrac{u\,\mathrm{d}u}{u} = \int \mathrm{d}u = u + C = \sqrt{x^2+9} + C$.  
Vì $F(0) = 5 \iff \sqrt{0+9} + C = 5 \iff 3 + C = 5 \iff C = 2$.  
Vậy $F(x) = \sqrt{x^2+9} + 2$.  
Khi đó $F(4) = \sqrt{4^2+9} + 2 = \sqrt{25} + 2 = 5 + 2 = 7$.

\item \textbf{Đáp số: $-1$.}  
Tính $\int (2x - 1)\mathrm{e}^x\,\mathrm{d}x$.  
Đặt $u = 2x - 1 \implies \mathrm{d}u = 2\,\mathrm{d}x$; $\mathrm{d}v = \mathrm{e}^x\,\mathrm{d}x \implies v = \mathrm{e}^x$.  
Áp dụng công thức từng phần:  
$\int (2x-1)\mathrm{e}^x\,\mathrm{d}x = (2x-1)\mathrm{e}^x - \int 2\mathrm{e}^x\,\mathrm{d}x = (2x-1)\mathrm{e}^x - 2\mathrm{e}^x + C = (2x - 3)\mathrm{e}^x + C$.  
Do đó $a = 2, b = -3 \implies T = a + b = 2 + (-3) = -1$.

\item \textbf{Đáp số: $2$.}  
Tính $\int_0^1 (2x+1)\mathrm{e}^x\,\mathrm{d}x$.  
Đặt $u = 2x+1 \implies \mathrm{d}u = 2\,\mathrm{d}x$; $\mathrm{d}v = \mathrm{e}^x\,\mathrm{d}x \implies v = \mathrm{e}^x$.  
$\int_0^1 (2x+1)\mathrm{e}^x\,\mathrm{d}x = [(2x+1)\mathrm{e}^x]_0^1 - 2\int_0^1 \mathrm{e}^x\,\mathrm{d}x = (3\mathrm{e} - 1) - 2[\mathrm{e}^x]_0^1$  
$= 3\mathrm{e} - 1 - 2(\mathrm{e} - 1) = \mathrm{e} + 1$.  
Suy ra $a = 1, b = 1 \implies P = a^2 + b^2 = 1^2 + 1^2 = 2$.

\item \textbf{Đáp số: $5$.}  
$I = \int_2^3 \dfrac{1}{x^2-1}\,\mathrm{d}x = \dfrac{1}{2}\int_2^3 \left(\dfrac{1}{x-1} - \dfrac{1}{x+1}\right)\,\mathrm{d}x = \dfrac{1}{2}\left[ \ln\left|\dfrac{x-1}{x+1}\right| \right]_2^3$.  
Tại $x = 3$: $\ln\left(\dfrac{2}{4}\right) = \ln\left(\dfrac{1}{2}\right) = -\ln 2$.  
Tại $x = 2$: $\ln\left(\dfrac{1}{3}\right) = -\ln 3$.  
Do đó $I = \dfrac{1}{2}\left[ -\ln 2 - (-\ln 3) \right] = \dfrac{1}{2}(\ln 3 - \ln 2) = \dfrac{1}{2}\ln\left(\dfrac{3}{2}\right)$.  
Suy ra $a = 3, b = 2 \implies a + b = 3 + 2 = 5$.

\item \textbf{Đáp số: $11$.}  
Đặt $t = \sqrt{x+1} \implies t^2 = x+1 \implies x = t^2 - 1 \implies \mathrm{d}x = 2t\,\mathrm{d}t$.  
Đổi cận: $x = 0 \implies t = 1$; $x = 3 \implies t = 2$.  
$I = \int_1^2 \dfrac{t^2-1}{t} \cdot 2t\,\mathrm{d}t = 2\int_1^2 (t^2 - 1)\,\mathrm{d}t = 2\left[ \dfrac{t^3}{3} - t \right]_1^2$  
$= 2\left[ \left(\dfrac{8}{3} - 2\right) - \left(\dfrac{1}{3} - 1\right) \right] = 2\left[ \dfrac{2}{3} - \left(-\dfrac{2}{3}\right) \right] = 2 \cdot \dfrac{4}{3} = \dfrac{8}{3}$.  
Suy ra $a = 8, b = 3 \implies a + b = 8 + 3 = 11$.

\item \textbf{Đáp số: $3$.}  
Nhận xét rằng $(x^2+x)' = 2x+1$.  
Do đó $\int_1^2 \dfrac{2x+1}{x^2+x}\,\mathrm{d}x = \int_1^2 \dfrac{\mathrm{d}(x^2+x)}{x^2+x} = \left[ \ln|x^2+x| \right]_1^2 = \ln 6 - \ln 2 = \ln\left(\dfrac{6}{2}\right) = \ln 3$.  
Vậy $k = 3$.

\item \textbf{Đáp số: $12$.}  
Đặt $u = \ln x \implies \mathrm{d}u = \dfrac{1}{x}\,\mathrm{d}x$; $\mathrm{d}v = x^2\,\mathrm{d}x \implies v = \dfrac{x^3}{3}$.  
$\int_1^{\mathrm{e}} x^2 \ln x\,\mathrm{d}x = \left[ \dfrac{x^3}{3}\ln x \right]_1^{\mathrm{e}} - \dfrac{1}{3}\int_1^{\mathrm{e}} x^2\,\mathrm{d}x = \left(\dfrac{\mathrm{e}^3}{3} - 0\right) - \dfrac{1}{3}\left[ \dfrac{x^3}{3} \right]_1^{\mathrm{e}}$  
$= \dfrac{\mathrm{e}^3}{3} - \dfrac{1}{9}(\mathrm{e}^3 - 1) = \dfrac{3\mathrm{e}^3 - \mathrm{e}^3 + 1}{9} = \dfrac{2\mathrm{e}^3 + 1}{9}$.  
Suy ra $a = 2, b = 1, c = 9 \implies S = a + b + c = 2 + 1 + 9 = 12$.

\item \textbf{Đáp số: $0{,}25$.}  
Đặt $t = \sin x \implies \mathrm{d}t = \cos x\,\mathrm{d}x$.  
Đổi cận: $x = 0 \implies t = 0$; $x = \dfrac{\pi}{2} \implies t = 1$.  
$I = \int_0^1 t^3\,\mathrm{d}t = \left[ \dfrac{t^4}{4} \right]_0^1 = \dfrac{1}{4} = 0{,}25$.

\item \textbf{Đáp số: $4$.}  
Ta có: $\int_0^m (2x - 3)\,\mathrm{d}x = \left[ x^2 - 3x \right]_0^m = m^2 - 3m$.  
Theo đề bài: $m^2 - 3m = 4 \iff m^2 - 3m - 4 = 0 \iff (m - 4)(m + 1) = 0$.  
Vì $m > 0$ nên $m = 4$.

\item \textbf{Đáp số: $4$.}  
Đặt $t = 3x \implies \mathrm{d}t = 3\,\mathrm{d}x \implies \mathrm{d}x = \dfrac{1}{3}\mathrm{d}t$.  
Đổi cận: $x = 0 \implies t = 0$; $x = 2 \implies t = 6$.  
Khi đó: $J = \int_0^2 f(3x)\,\mathrm{d}x = \dfrac{1}{3}\int_0^6 f(t)\,\mathrm{d}t = \dfrac{1}{3} \cdot 12 = 4$.

\item \textbf{Đáp số: $50$.}  
Ô tô dừng lại hẳn khi vận tốc $v(t) = 0 \iff 20 - 4t = 0 \iff t = 5\text{ s}$.  
Quãng đường ô tô di chuyển được là:  
$s = \int_0^5 (20 - 4t)\,\mathrm{d}t = \left[ 20t - 2t^2 \right]_0^5 = 20(5) - 2(5^2) = 100 - 50 = 50\text{ m}$.

\item \textbf{Đáp số: $35$.}  
Phương trình hoành độ giao điểm của parabol với $Ox$: $-x^2 + 4x = 0 \iff x = 0$ hoặc $x = 4$.  
Trên đoạn $[0; 4]$, $-x^2 + 4x \ge 0$.  
Diện tích hình phẳng:  
$S = \int_0^4 (-x^2 + 4x)\,\mathrm{d}x = \left[ -\dfrac{x^3}{3} + 2x^2 \right]_0^4 = -\dfrac{64}{3} + 32 = \dfrac{32}{3}$.  
Suy ra $a = 32, b = 3 \implies a + b = 32 + 3 = 35$.

\item \textbf{Đáp số: $4{,}5$.}  
Phương trình hoành độ giao điểm: $x^2 - 2x = x \iff x^2 - 3x = 0 \iff x = 0$ hoặc $x = 3$.  
Trên đoạn $[0; 3]$, $x \ge x^2 - 2x$.  
Diện tích hình phẳng là:  
$S = \int_0^3 (x - x^2 + 2x)\,\mathrm{d}x = \int_0^3 (3x - x^2)\,\mathrm{d}x = \left[ \dfrac{3x^2}{2} - \dfrac{x^3}{3} \right]_0^3 = \dfrac{27}{2} - 9 = \dfrac{9}{2} = 4{,}5$.

\item \textbf{Đáp số: $16$.}  
Chọn hệ trục tọa độ $Oxy$ sao cho gốc $O$ là trung điểm của cạnh đáy dưới hình chữ nhật dài $6\text{ m}$, trục $Ox$ trùng với cạnh đáy này. Khi đó cạnh đáy nằm trên đoạn $[-3; 3]$ trên trục hoành.  
Đỉnh của parabol nằm tại trung điểm của cạnh trên, cách cạnh đáy $4\text{ m}$, nên đỉnh có tọa độ $I(0; 4)$.  
Parabol đi qua hai điểm góc phía dưới là $(-3; 0)$ và $(3; 0)$.  
Phương trình parabol có dạng $y = ax^2 + 4$.  
Thay điểm $(3; 0)$ vào ta được: $a(3)^2 + 4 = 0 \iff 9a + 4 = 0 \iff a = -\dfrac{4}{9}$.  
Vậy phương trình parabol là $y = 4 - \dfrac{4}{9}x^2$.  
Diện tích khu vực trồng hoa hồng:  
$S = \int_{-3}^3 \left(4 - \dfrac{4}{9}x^2\right)\,\mathrm{d}x = 2\int_0^3 \left(4 - \dfrac{4}{9}x^2\right)\,\mathrm{d}x = 2\left[ 4x - \dfrac{4}{27}x^3 \right]_0^3 = 2(12 - 4) = 16\text{ m}^2$.

\item \textbf{Đáp số: $9$.}  
Thể tích khối tròn xoay thu được khi quay hình phẳng quanh $Ox$:  
$V = \pi \int_0^3 (\sqrt{2x})^2\,\mathrm{d}x = \pi \int_0^3 2x\,\mathrm{d}x = \pi \left[ x^2 \right]_0^3 = 9\pi$.  
Vậy $\dfrac{V}{\pi} = 9$.

\item \textbf{Đáp số: $37$.}  
Thể tích khối tròn xoay:  
$V = \pi \int_0^2 (x^2)^2\,\mathrm{d}x = \pi \int_0^2 x^4\,\mathrm{d}x = \pi \left[ \dfrac{x^5}{5} \right]_0^2 = \dfrac{32\pi}{5}$.  
Suy ra $\dfrac{a}{b} = \dfrac{32}{5} \implies a = 32, b = 5 \implies a + b = 32 + 5 = 37$.

\item \textbf{Đáp số: $40$.}  
Thể tích chứa nước của chiếc bình:  
$V = \pi \int_0^8 (\sqrt{x+1})^2\,\mathrm{d}x = \pi \int_0^8 (x+1)\,\mathrm{d}x = \pi \left[ \dfrac{x^2}{2} + x \right]_0^8 = \pi \left(\dfrac{64}{2} + 8\right) = 40\pi\text{ dm}^3\text{ (lít)}$.  
Vậy $\dfrac{V}{\pi} = 40$.

\item \textbf{Đáp số: $35$.}  
Mặt cắt tại điểm $x$ là hình vuông có cạnh bằng $\sqrt{4x - x^2}$ nên diện tích thiết diện là:  
$S(x) = (\sqrt{4x - x^2})^2 = 4x - x^2$.  
Thể tích của vật thể là:  
$V = \int_0^4 S(x)\,\mathrm{d}x = \int_0^4 (4x - x^2)\,\mathrm{d}x = \left[ 2x^2 - \dfrac{x^3}{3} \right]_0^4 = 2(16) - \dfrac{64}{3} = 32 - \dfrac{64}{3} = \dfrac{32}{3}$.  
Suy ra $a = 32, b = 3 \implies a + b = 32 + 3 = 35$.
\end{enumerate}

\end{document}
